\par\bigskip\Needspace{6\baselineskip}
\pdfbookmark[0]{Bemerkungen zur Analysis residuorum}{bemerkungen-analysis-residuorum}
\begin{center}\large\bfseries BEMERKUNGEN ZUR ANALYSIS RESIDUORUM.\end{center}
\BookRule{.09\linewidth}
\renewcommand{\BookRunningTitle}{Nachlass. Bemerkungen zur Analysis residuorum}

Die beiden vorstehenden Abhandlungen sind einem umfangreichen Manu\-scripte entnommen, welches den Titel Analysis Residuorum führt und vermuthlich aus dem Jahre 1797 oder 1798 stammt; durch eine gänzliche Umarbeitung sind aus demselben später die Disquisitiones Arithmeticae entstanden. Der vollständige Titel des Caput sextum lautet:

Solutio congruentiae $x^m-1\equiv0$ et aequationis $x^m-1=0$; cum dilucidationibus super theoria polygonorum regularium.

Der zweite Theil desselben (\textsection\textsection.\ 253--278) ist seinem wesentlichen Inhalte nach in die siebente Section der Disqq. Arithm. übergegangen.

Ausserdem ist noch zum Theil erhalten das Caput septimum. Variae quarundam investigationum praecedentium applicationes (\textsection\textsection.\ 279--302). Es zerfällt in folgende Unterabtheilungen:

De fractionum communium transmutationibus (\textsection\textsection.\ 279--281).

De fractionum communium in decimales conversione (\textsection\textsection.\ 282--292).

De resolutione aequationis indeterminatae $xx=a+by$ (\textsection\textsection.\ 293--297).

De resolutione aequationis indeterminatae $axx+byy=c$ (\textsection\textsection.\ 298--301).

De investigatione divisorum numerorum (\textsection.\ 302; die folgenden Bogen fehlen).

Dies alles ist fast wörtlich in die sechste Section der Disqq. Arithm. aufgenommen.

Die beiden hier mitgetheilten Abschnitte behandeln die Gegenstände, welche, wie aus der Vorrede und den Artikeln 11, 44, 61, 62, 65, 84 der Disqq. Arithm. hervorgeht, den Inhalt der achten Section dieses Werkes bilden sollten. Es verdient indessen bemerkt zu werden, dass dieser Plan später wieder abgeändert ist; es findet sich nemlich unter den Manuscripten ein Fragment mit der Ueberschrift Sectio octava: Quarundam disquisitionum ad circuli sectionem pertinentium uberior consideratio. Dasselbe 
\SourceLeaf{251}
beginnt mit Art.\ 367 und sollte also die Fortsetzung der Disqq.\ Arithm.\ bilden; die wenigen noch vorhandenen Artikel sind aber später ihrem Inhalte nach in die Abhandlung Summatio quarundam serierum singularium übergegangen, und deshalb wird dieses Fragment von der gegenwärtigen Ausgabe ausgeschlossen.

In dem vorstehenden Abdruck der beiden Theile der Analysis Residuorum ist der Text des Originals im Wesentlichen treu beibehalten, obgleich dasselbe in formeller Beziehung nicht druckfertig zu nennen ist; in den folgenden Bemerkungen sind die wichtigsten Abänderungen bezeichnet, und zugleich einige Erläuterungen hinzugefügt.

\textsection.\ 237.\ Vergl.\ Disqq.\ Arithm.\ artt.\ 61, 62.

\textsection.\ 239.\ Vergl.\ Disqq.\ Arithm.\ artt.\ 53, 54, 65.

\textsection.\ 241.\ Wenn $n=2^\nu$ und $\nu\geqq3$ ist, so existirt zwar keine Zahl $\rho$ von der angegebenen Art, aber die ganze Untersuchung wird hierdurch nicht wesentlich geändert.

\textsection.\ 251.\ Vermuthlich sollte die hier bemerkte Schwierigkeit durch die Einführung höherer Potenzen von $p$ als Moduln beseitigt werden. Vergl.\ \textsection\textsection.\ 363, 372, 373.

\textsection.\ 332.\ Die Voraussetzung, dass der Modulus eine Primzahl ist, wird bis \textsection.\ 372 incl.\ beibehalten.

\textsection.\ 338.\ Das unvollständige Citat kann auf Disqq.\ Arithm.\ art.\ 44 bezogen werden.

\textsection\textsection.\ 344--346.\ Von den beiden im Manuscript vorhandenen Beweisen ist hier der erste, welcher mit den Worten iam demonstrare accingimur eingeleitet wird und sich auf eine nähere Untersuchung der Ausdrücke $(1^\alpha2^\beta3^\gamma\ldots)$ gründet, nach der eigenen Vorschrift des Verfassers ganz unterdrückt (`Tota praecedens demonstratio una cum altera theorematis praec., quam adiicere mens erat, supprimenda erit, quoniam aliam infinities simpliciorem deteximus. Nititur ea huic fundamento'); in dem obigen Abdruck ist ferner der zweite Beweis dadurch abgekürzt, dass die Entwicklung von $\frac{x\,\mathrm{d}P}{P\,\mathrm{d}x}$ statt derjenigen von $\frac{x\,\mathrm{d}P}{\mathrm{d}x}$ betrachtet wird, wodurch zugleich eine im Original enthaltene Beziehung auf den unterdrückten ersten Beweis umgangen wird.

\textsection.\ 348.\ Der Ausdruck radix prima ist hier in derselben Bedeutung zu nehmen, wie der Ausdruck radix propria in der Abhandlung Summatio quarundam serierum singularum art.\ 11. --- Bei der Behauptung, dass die Coëfficienten $A'.B'\ldots$ des entwickelten Productes ganze rationale Zahlen sind, wird auf das sechste Capitel verwiesen, in welchem aber die Theorie der Gleichung $x^\tau-1=0$ nur für den Fall behandelt wird, dass $\tau$ eine Primzahl ist; die Form des Beweises in \textsection.\ 349 führt zunächst auf folgende Ergänzung. Wird das entwickelte Product in die (für alle Wurzeln der Gleichung $\theta^\tau=1$ geltende) Form
\[
S=E+F\theta+\ldots+N\theta^{\tau-1}
\]
gebracht, so sind die Coëfficienten $E,F\ldots N$ ganze rationale Functionen von $x$ mit ganzen rationalen Coëfficienten; da ferner das Product ungeändert bleibt, wenn $\theta$ durch $\theta^k$ ersetzt wird, wo $k$ irgend eine relative Primzahl zu $\tau$ bedeutet, so gilt dasselbe von dem Ausdruck $S$, und hieraus ergiebt sich ohne Schwierigkeit, dass alle diejenigen in $S$ enthaltenen Potenzen von $\theta$, deren Exponenten $s$ einen und denselben grössten gemeinschaftlichen Divisor mit $\tau$ haben, auch identische Coëfficienten haben müssen; da endlich eine jede Summe solcher Potenzen $\theta^s$ immer eine ganze Zahl ist, so leuchtet ein, dass der Ausdruck $S$, und folglich auch das in Rede stehende Product eine ganze Function von $x$ mit ganzen Coëfficienten ist, was zu zeigen war. Ebenso geht aus dieser Betrachtung zugleich die Richtigkeit der Bemerkung am Schlusse des Paragraphen hervor. Andere Gründe lassen indessen vermuthen, dass dem Verfasser schon damals das allgemeine Theorem über die Transformation der symmetrischen Functionen (Demonstratio nova altera theorematis omnem functionem etc.\ art.\ 4) bekannt war, aus welchem sich die obigen Sätze als unmittelbare Folgerungen ergeben.

\textsection.\ 352.\ Das Zeichen $R\equiv S(\mathrm{mod.}\,P)$ oder auch $R\equiv S(\mathrm{mod.}\,P,p)$ bedeutet hier und im 
\SourceLeaf{252}
Folgenden, dass die Differenz $R-S$ nach dem Modul $p$ den Divisor $P$ hat. --- Das unvollständige Citat kann auf Disqq.\ Arithm.\ art.\ 49 bezogen werden.

\textsection.\ 354.\ Durch Multiplication mit $x^\nu-1$ ergiebt sich, dass die Summen gleich hoher Potenzen der Wurzeln der beiden Gleichungen $(P,\rho^{k\nu+t})=0$, $(P,\rho^t)=0$ einander congruent sind $(\mathrm{mod.}\,p)$, und hieraus folgt die Congruenz $(P,\rho^{k\nu+t})\equiv(P,\rho^t)(\mathrm{mod.}\,p)$, sobald $m<p$ ist (vergl.\ \textsection.\ 244); ist aber $m\geqq p$, so lässt sich der Coëfficient der Potenz $x^{m-p}$ in einer Gleichung nicht mehr aus den gegebenen Potenzsummen ihrer Wurzeln nach dem Modul $p$ bestimmen, weil er in den hierzu dienenden \textsc{Newton}'schen Formeln mit dem Factor $p$ behaftet ist. In der That darf man aus der Congruenz je zweier gleich hoher Potenzsummen der Wurzeln der Gleichungen $A=0$, $B=0$ allgemein nur folgern, dass $A\equiv\mathfrak A^p\mathfrak C$, $B=\mathfrak B^p\mathfrak C(\mathrm{mod.}\,p)$ ist, wo $\mathfrak C$ den grössten gemeinschaftlichen Divisor der beiden Functionen $A,B$ nach dem Primzahl-Modulus $p$ bezeichnet, $\mathfrak A$ und $\mathfrak B$ aber ganz unbestimmte Functionen sind. Es ist zu vermuthen, dass der Verfasser die Allgemeingültigkeit des Satzes aus der Theorie der Transformation der symmetrischen Functionen und speciell aus dem folgenden Satze abgeleitet hat: Ist in Bezug auf einen beliebigen Modulus $p$ die Differenz $R(x)-S(x)$ theilbar durch die Function $P(x)$, und sind $a,b,c\ldots$ die Wurzeln der Gleichung $P(x)=0$, so sind die Functionen
\[
\begin{gathered}
(x-R(a))(x-R(b))(x-R(c))\ldots\quad\text{und}\\
(x-S(a))(x-S(b))(x-S(c))\ldots
\end{gathered}
\]
einander nach dem Modul $p$ congruent.

\textsection.\ 355.\ Es wird in \textsection.\ 368 gezeigt, dass $P$ und $\frac{\mathrm{d}P}{\mathrm{d}x}$ keinen gemeinschaftlichen Divisor haben, wenn $P$ keinen Factor mehr als einmal enthält.

\textsection\textsection.\ 358, 359.\ Die unter den Text gesetzte Note ist einem einzelnen Blatt entnommen, welches wahrscheinlich den schon in der Handschrift gestrichenen \textsection.\ 359 ersetzen sollte.

\textsection.\ 360.\ In dem Ausdruck des Theorems ist eine Ungenauigkeit der Handschrift berichtigt.

\textsection.\ 361.\ Hier bedeutet der Exponent $\frac1k$ in dem Zeichen $(\xi,\rho^{\frac1k})$ jede positive ganze Zahl $k'$ von der Beschaffenheit, dass $kk'\equiv1(\mathrm{mod.}\,\nu)$ wird, wo $\nu$ die kleinste positive ganze Zahl ist, für welche $x^\nu-1$ durch $\xi$ nach dem Modul $p$ theilbar wird; hierbei ist vorauszusetzen, dass $\xi$ nicht durch $x$ theilbar nach dem Modul $p$, und ausserdem, dass $k$ relative Primzahl zu $\nu$ ist. Die Richtigkeit der Behauptung, dass $\xi'$ durch $(\xi,\rho^{\frac1k})$ theilbar ist $(\mathrm{mod.}\,p)$, ergiebt sich aus \textsection.\ 354.

\textsection.\ 363.\ Die Schlussbemerkung bezieht sich vermuthlich auf die Einführung von Moduln, welche Potenzen der Primzahl $p$ sind; vergl.\ \textsection\textsection.\ 251, 372, 373.

\textsection.\ 367.\ Die Wurzeln der Gleichung $x^3+xx-2x-1=0$ sind die zweigliedrigen Perioden, in welche die Wurzeln der Gleichung $\frac{x^7-1}{x-1}=0$ zerfallen. Dasselbe Beispiel findet sich auch auf einem einzelnen Blatt, wo das Hauptresultat der \textsection\textsection.\ 362, 363 unter dem Titel `der goldene Lehrsatz' ausgesprochen ist.

\textsection.\ 371.\ Dieser Paragraph sollte ein Beispiel enthalten; doch ist dasselbe nicht ausgeführt.
\begin{flushright}R.\ \textsc{Dedekind}.\end{flushright}
\BookRule{.28\linewidth}
